\subsection{\texorpdfstring{SUMMABLE FAMILIES AND MULTIPLIABLE FAMILIES IN \(\overline{\mathbf{R}}\)}{SUMMABLE FAMILIES AND MULTIPLIABLE FAMILIES IN \textbackslash overline\{\textbackslash mathbf\{R\}\}}}

On the interval \([0, +\infty]\) of \(\overline{\mathbf{R}}\) , addition is an associative and commutative law of composition (§ 4, no. 3); hence the notion of a summable family of numbers in this interval is defined (Chapter III, § 5, no. 1, Remark 3).

PROPOSITION 2. Every family \((x_{\mathfrak{t}})\) of real numbers \(\geqslant 0\) is summable in \(\overline{\mathbf{R}}\) . For the mapping \(\mathbf{H} \to s_{\mathbf{H}}\) of the directed set \(\mathfrak{F}(\mathbf{I})\) into \(\overline{\mathbf{R}}\) is increasing; hence (§ 5, no. 2, Theorem 2) has a limit.

The same reasoning shows that every family of real numbers \(\leqslant 0\) is summable in \(\overline{\mathbf{R}}\) .

Similarly, multiplication is an associative and commutative law of composition on each of the intervals \([0, 1]\) and \([1, +\infty]\) of \(\overline{R}\) , and therefore the notion of a multipliable family is defined on each of these intervals.

PROPOSITION 3. Every family \((1 + u_{t})\) {[}resp. \((1 - u_{t})]\) of numbers \(\geqslant 1\) (resp. \(\geqslant 0\) and \(\leqslant 1\) ) is multipliable in \(\overline{\mathbf{R}}\) .

Same proof as for Proposition 2.

COROLLARY. The product \(\prod_{i} (1 + u_i) [\text{resp. } \prod_{i} (1 - u_i))\) of numbers \(\geqslant 1\) (resp. \(>0\) and \(\leqslant 1\) ) is equal to \(+\infty\) (resp. 0) if and only if \(\sum_{i} u_i = +\infty\) . For if \(\sum_{i} u_i\) is finite then \(\prod_{i} (1 + u_i)\) and \(\prod_{i} (1 - u_i)\) are in \(\mathbf{R}^*\) , and conversely (no. 4, Theorem 4).

Remark. The theorem of associativity (Chapter III, § 5, no. 3, Theorem 2) remains valid when G is replaced by \(\overline{\mathbf{R}}\) and the \(x_{i}\) are assumed to be \(\geqslant 0\) . This is clear if \(\sum_{i \in I} x_{i}\) is finite. Suppose, on the contrary, that \(\sum_{i \in I} x_{i} = +\infty\) . Then for each finite \(a > 0\) there is a finite subset H of I such that \(\sum_{i \in H} x_{i} \geqslant a\) . Let K be a finite subset of L such that \(H \subset \bigcup_{\lambda \in K} I_{\lambda}\) ; then since \(s_{\lambda} \geqslant \sum_{i \in I_{\lambda} \cap H} x_{i}\) for all \(\lambda \in K\) , we have

\[
\sum_ {\lambda \in \mathbf {K}} s _ {\lambda} \geqslant \sum_ {i \in \mathbf {H}} x _ {i} \geqslant a,
\]

which shows that \(\sum_{\lambda \in L} s_{\lambda} = +\infty\) . We leave to the reader the task of formulating the analogous proposition for multipliable families of numbers in \([0, 1]\) or in \([1, +\infty]\) .

